\section{Introduction}\label{sec:intro}

A new kind of contributor has joined software teams. Coding agents, language models given a
shell, a file system and a test runner, now read a repository, plan a change, edit many files
and run the tests before a person sees the result
\citep{chen2021codex,jimenez2024swebench,yang2024sweagent}. They are fast and broadly competent.
They are also, in one respect, the most difficult colleague a team has ever had: every session
starts with no memory of the last one. Whatever the agent learned about the domain yesterday,
about why the work-in-progress limit is checked in the service rather than the route, or why the
domain objects never import the web framework, is gone this morning.

The only thing a person and an agent reliably share is the repository. Whatever the repository
does not say, and does not enforce, the agent will reinvent, usually plausibly and often
differently each time. People have the opposite problem. They remember the domain, but they
cannot read every line an agent writes. Studies of programmers working with code-generating
tools describe a shift from writing to reviewing, and reviewing under that volume slides
towards skimming \citep{vaithilingam2022expectation,barke2023grounded}. Security studies show
what can pass unnoticed \citep{pearce2022asleep,perry2023insecure}.

We give the two resulting forces names. \emph{Sprawl} is what happens when a large volume of
locally plausible changes spreads one domain rule across many places. A check is copied into a
route because it was nearby. A notification is called inline because that was quicker than
publishing an event. A framework annotation lands on a domain object because it was convenient.
Every one of those changes would pass review on its own. Taken together they leave the domain
in pieces. \emph{Cognitive erosion} is what happens when neither party can hold enough of the
domain in mind to notice: the reviewer checks syntax because the intent is not visible in the
diff, and the agent re-derives conventions from scratch every session. Lehman observed long ago
that a system's structure decays under change unless work is done to preserve it
\citep{lehman1980laws}. Agents have not repealed that law. They have made it run faster.

\paragraph{Where we meet the problem} This paper comes out of a collaboration between a
university and an industrial engineering site in the same city, and we meet the problem from
both sides. At South East Technological University (SETU), students on the online Higher
Diploma in Science in Computer Science\footnote{\url{https://www.setu.ie/courses/higher-diploma-in-computer-science-2-years-online}} now learn to
program alongside coding agents, and their lecturers need codebases that teach a domain to both.
SETU's Master of Business Studies in Lean Enterprise
Excellence\footnote{\url{https://www.setu.ie/courses/master-of-business-in-lean-enterprise-excellence}} teaches the lean practice this
paper draws on. SETU delivers course material through Tutors, an open-source learning
platform whose production TypeScript codebase coding agents were already changing
\citep{tutorsmonorepo}. At Red Hat, the same agents contribute to engineering work, and hiring
is moving towards graduates who can work safely in repositories that agents also change. The
question is the same for a student, a new hire and an agent: how does a repository make its
domain rules impossible to miss and hard to break, without burying everyone in prose?

\paragraph{The need for constraints} The common response is to write more guidance: longer
prompts, longer contributor guides, longer instruction files for agents
\citep{agentsmd}. Guidance helps, but we have watched it fail in two ways. It is read at the
start of a session and then competes with everything else in the agent's context, and it grows
until nobody, person or agent, reads it at all. Our own first step beyond guidance was a
repository in which conventions are carried three ways at once: exemplified by one small
vertical slice, declared in a manifest of a few hundred words, and enforced by architecture
rules written as tests. The receivers in this paper began from that arrangement. This paper
generalises it. It asks not only \emph{whether} a rule is written down but \emph{how strongly}
it is asserted, and it extends the question from structure to behaviour and requirements.

\paragraph{The Assertion Ladder} Our central claim is a teaching one. A coding agent is a
capable learner with no memory, and the repository is its teacher. Teachers do not assert
everything the same way. They show a worked example, they tell a rule, they set an exercise that
marks the answer, they give feedback that explains the mistake, and now and then they arrange
the task so that the wrong answer is simply not available. The \emph{Assertion Ladder}
(Figure~\ref{fig:ladder}) gives each of these moves a rung, from \rung{L0 Implied} to
\rung{L5 Prevented}, and asks whoever designs a repository to place every domain rule on a rung
on purpose: high enough that the learner cannot miss it, and no higher than the rule's risk
justifies, because every assertion costs something to write, to maintain and to read.

\paragraph{Lean as the explanation} None of this is new engineering. It is old engineering
applied to a new worker. Lean thinking starts from \emph{value}, what the customer and, here,
the domain need, and treats everything that does not serve it as waste
\citep{womack1996lean,poppendieck2006implementing}. For work that repeats, its central tool is
\emph{standard work}: the current best method, followed by everyone who does the job and
improved by the team that owns it \citep{liker2004toyota}. It reads a process as a \emph{value
stream} and asks where along it a defect is found and what the waiting and rework cost
\citep{rother1999learning}. And it answers each problem with \emph{root cause analysis}, asking
why until the answer is something the process can change \citep{ohno1988tps}. Read this way,
the rungs are standard work for a contributor who cannot remember it: shown in a worked slice,
told in a short manifest, and then built into the work as checks that stop the build, explain
the stop, or make the mistake impossible to write. Section~\ref{sec:discussion} makes the case
that lean explains why the ladder works, rather than decorating it.

\paragraph{Critical gaps} Software engineering already has most of the parts. Architecture
conformance checking, specification by example, structured requirements and instruction files
for agents all exist, and Section~\ref{sec:related} reviews them. What it lacks is a standard
way to combine them. We identify five gaps in current practice that call for standardisation.
There is no shared vocabulary for how strongly a rule is asserted (\ref{gap:vocabulary}).
Instruction files for agents are unvalidated prose with no link to enforcement
(\ref{gap:prose}). Conformance tools say nothing about which rules to check rather than
document, or about a failure naming its fix (\ref{gap:selection}). Agent-facing repositories
have no executable closure between requirements and the scenarios that test them
(\ref{gap:closure}). And there is no benchmark of whether a repository's own harness catches
the mistakes that matter, as opposed to whether an agent can complete a task
(\ref{gap:benchmark}). None of these gaps is large on its own. Their combined effect is that
every team sharing a repository with agents invents its own answer, and most answer with more
prose.

\paragraph{Two known uses} We show the pattern in two settings. The first is greenfield:
one application, \emph{taskboard}, a small Kanban service with a work-in-progress limit, built
with the ladder from the start. We walk through its TypeScript implementation and show that
Java and Python implementations with identical behaviour bind every rung equivalently
\citep{taskboardrepo}. We call these implementations \emph{receivers} of the pattern rather
than its source, and work backwards from what each gains to the single capability they share.
The second is a retrofit: \emph{Tutors}, an open-source TypeScript learning platform in
production, to which the rungs were added after the fact while coding agents were already
contributing, and from which the release checks were then extracted into a reusable,
product-independent harness \citep{tutorsmonorepo,tutorsreleaseharness}. One shows the ladder
built where nothing existed. The other shows it fitted to a codebase with years of history.

\paragraph{Contributions} The paper makes five contributions, each aimed at one or more of
the gaps.
\begin{enumerate}
  \item The Assertion Ladder, a pattern in the classic form \citep{alexander1977pattern,gamma1994design}
    with six rungs and nine rung patterns, proposed as a shared vocabulary for how strongly a
    repository asserts each domain rule (\ref{gap:vocabulary}, \ref{gap:prose},
    \ref{gap:selection}; Section~\ref{sec:ladder}).
  \item A unification of three kinds of intent, structure, behaviour and requirements, as three
    loads on the same ladder, with a closure check that ties requirements to scenarios and the
    observation that the pattern of an EARS requirement \citep{mavin2009ears} predicts the layer
    that should implement it (\ref{gap:closure}; Sections~\ref{sec:ladder}
    and~\ref{sec:receivers}).
  \item A greenfield receiver in TypeScript, with equivalent Java and Python receivers, showing
    which rungs each language can reach and what it uses where it cannot
    (Section~\ref{sec:receivers}).
  \item A retrofit of a production codebase, Tutors, and the extraction of its release checks
    into a reusable harness, showing how the ladder is adopted where it was not designed in
    (Section~\ref{sec:tutors}).
  \item Violation seeding as a benchmark of a repository rather than of an agent, run before and
    after the behavioural and requirement rungs were added, and a pilot with eighteen coding
    agents (\ref{gap:benchmark}; Section~\ref{sec:evaluation}).
\end{enumerate}

\begin{figure*}[t]
  \centering
  \resizebox{\linewidth}{!}{% The Assertion Ladder: six rungs, each with its teaching move and its lean counterpart.
\begin{tikzpicture}[
    font=\small,
    rung/.style={draw=#1, line width=0.9pt, fill=#1!8, rounded corners=2pt,
                 minimum width=5.9cm, minimum height=0.78cm, anchor=west, align=left},
    note/.style={anchor=west, align=left, text width=4.6cm, font=\footnotesize},
  ]
  \def\gap{0.95}
  \foreach \i/\name/\teach/\lean/\col in {
      0/{L0 Implied}/{hidden curriculum}/{tacit knowledge, no standard}/muted,
      1/{L1 Shown}/{worked example}/{standard work, shown}/structure,
      2/{L2 Told}/{direct instruction}/{standard work, written}/requirement,
      3/{L3 Checked}/{assessment}/{built-in test at the point of work}/behaviour,
      4/{L4 Explained}/{formative feedback}/{andon: the stop says where and why}/reqc,
      5/{L5 Prevented}/{constraint}/{mistake-proofing (poka-yoke)}/structure} {
    \node[rung=\col] (r\i) at (0, \i*\gap) {\textbf{\name}\quad\textcolor{muted}{\footnotesize\teach}};
    \node[note] at (6.2, \i*\gap) {\textcolor{muted}{lean:} \lean};
  }
  % the rails of the ladder
  \draw[line width=1.6pt, jotblue] (-0.25,-0.5) -- (-0.25, 5*\gap+0.5);
  % arrows
  \draw[-{Stealth[length=3mm]}, line width=1pt, jotblue] (-0.8,-0.3) -- node[rotate=90, above, font=\footnotesize] {earlier, more certain feedback} (-0.8, 5*\gap+0.3);
  \draw[-{Stealth[length=3mm]}, line width=1pt, muted] (11.1, 5*\gap+0.3) -- node[rotate=-90, above, font=\footnotesize] {cheaper to build and maintain} (11.1,-0.3);
\end{tikzpicture}
}
  \caption{The Assertion Ladder. Each rung is a teaching move and a lean mechanism. Rules climb
  only as far as their risk warrants; as a rule reaches a checked rung, it is removed from the
  prose below it.}
  \label{fig:ladder}
\end{figure*}
